% ===== 1 问题定义与贡献 =====
\section{问题定义与主要贡献}
\label{sec:r-problem}

骨料粒径级配与颗粒形貌直接影响混凝土的堆积结构、工作性能和力学性能。工程检测通常依据 \texttt{GB/T 14685}、\texttt{JGJ 52} 等标准，通过逐级机械筛分和称重获得筛级质量占比及累计通过率\cite{gb14685-2022,jgj52-2006}。其物理意义明确，但取样、振筛、称量和记录均需要离线操作，难以直接嵌入生产过程形成快速反馈。本赛题因此要求参赛系统在自然堆积或输送场景中，以非人工筛分方式完成颗粒精细识别、粒径分析、形貌识别和异常检测，并在有限现场时间内给出可评价结果。

本任务与一般目标检测的根本区别在于：\textbf{比赛最终评价的是标准筛分语义，而不是图像中目标框或像素本身。} 从一幅自然堆积图像得到标准筛分级配，需要连续跨越三个层次。第一，自然堆积中的粘连、遮挡和阴影使“前景”必须进一步分解为独立颗粒；merge 会制造偏大的伪颗粒，split 会制造偏小实例。第二，像素轮廓必须通过本次成像条件的尺度标定转换为毫米量，否则长短轴和 50~mm 超限阈值没有物理意义。第三，也是最关键的一层，单目图像观察的是二维投影并天然按颗粒个数统计，而机械筛分反映三维颗粒的筛孔通过行为并按质量称重。因此，即使实例边界和毫米尺寸均准确，直接把二维短轴或数量分布当作机械筛分结果仍可能存在系统偏差。

图\ref{fig:r-three-gaps}把这三个问题组织为一条误差传播链。它们不是彼此独立的三个“功能点”：上游实例错误会改变几何量，几何量再经过质量幂次放大后影响场景级分位数。例如，两个相邻颗粒若被 merge 为一个实例，其短轴和面积通常增大；该伪粗颗粒又会在 $d^\gamma$ 权重下获得更大的质量贡献，因此最终偏差可能明显大于普通像素误差本身。相反，一个颗粒被 split 后会形成多个偏细实例，使数量统计和低粒径端同时增加。由此，本文不把实例分割精度、毫米尺度和级配映射分别优化后简单相加，而是始终围绕最终筛分语义检查误差如何向下游传播。

\begin{figure}[t]
\centering
\begin{tikzpicture}[
  node distance=0.55cm,
  stage/.style={rectangle,draw,rounded corners,minimum width=4.0cm,minimum height=1.0cm,align=center,font=\footnotesize},
  err/.style={rectangle,draw,dashed,rounded corners,minimum width=3.8cm,minimum height=0.85cm,align=center,font=\scriptsize},
  arr/.style={->,thick}
]
\node[stage,fill=blue!7] (s1) {图像 $I$ $\rightarrow$ 独立实例 $\{M_i\}$\\\textbf{鸿沟 1：接触/遮挡颗粒如何分离}};
\node[stage,fill=green!7,right=0.65cm of s1] (s2) {实例 $\{M_i\}$ $\rightarrow$ 毫米几何 $\mathbf{x}_i$\\\textbf{鸿沟 2：pixel 如何获得物理尺度}};
\node[stage,fill=orange!8,right=0.65cm of s2] (s3) {二维几何 $\rightarrow$ 标准筛分 $\mathcal{Y}$\\\textbf{鸿沟 3：投影/计数如何对应筛孔/称重}};
\node[err,below=0.55cm of s1] (e1) {merge：偏粗；split：偏细\\漏检/遮挡：代表性下降};
\node[err,below=0.55cm of s2] (e2) {尺度偏差近似按比例移动\\全部粒径和筛级边界};
\node[err,below=0.55cm of s3] (e3) {二维 $\neq$ 三维筛孔行为\\数量分布 $\neq$ 质量分布};
\draw[arr] (s1)--(s2); \draw[arr] (s2)--(s3);
\draw[arr] (s1)--(e1); \draw[arr] (s2)--(e2); \draw[arr] (s3)--(e3);
\end{tikzpicture}
\caption{自然堆积视觉筛分的三个连续语义鸿沟及其典型误差方向。系统设计的目标不是单独“做好分割”，而是控制这些误差向最终标准筛分结果的传播。}
\label{fig:r-three-gaps}
\end{figure}

据此，本文把任务形式化为一条分层测量链。给定图像 $I$，实例分割得到 $N$ 个可见颗粒掩膜
\begin{equation}
\mathcal{M}=\{M_i\}_{i=1}^{N}.
\end{equation}
对每个实例，经尺度恢复获得几何描述 $\mathbf{x}_i=(a_i,b_i,A_i,P_i,C_i,S_i,\ldots)$，其中 $a_i,b_i$ 为投影长短轴，$A_i,P_i$ 为面积与周长，$C_i,S_i$ 为圆度和凸度。随后构造筛分等效粒径 $d_i=f(\mathbf{x}_i;\boldsymbol{\theta})$、形貌类别 $c_i$ 和异常状态 $z_i$；再以质量代理权重 $w_i=d_i^{\gamma}$ 聚合为质量相关累计分布，输出 $D_{10}/D_{50}/D_{90}$ 与各标准筛级质量占比。完整链条可概括为
\begin{equation}
I\rightarrow\{M_i\}\rightarrow\{\mathbf{x}_i\}\rightarrow\{d_i,w_i\}\rightarrow\mathcal{Y}_{\mathrm{sieve}}.
\label{eq:r-pipeline}
\end{equation}

表\ref{tab:r-task-map}给出赛题任务与系统证据的直接对应关系。这里刻意把“输出”和“验证证据”分开：系统能够生成某类结果，只能证明功能存在；只有使用与该任务匹配的 ground truth，才能进一步声明准确率。

\begin{table}[t]
\centering
\caption{赛题任务、系统输出与最终验证证据。}
\label{tab:r-task-map}
\footnotesize
\renewcommand{\arraystretch}{1.12}
\begin{tabular}{>{\raggedright\arraybackslash}p{2.0cm}>{\raggedright\arraybackslash}p{5.2cm}>{\raggedright\arraybackslash}p{6.2cm}}
\toprule
任务 & 系统输出 & 最终应使用的验证证据 \\
\midrule
精细识别 & 独立实例 mask、实例数量、长短轴 & 人工实例掩膜；IoU/F1、merge/split \\
粒径分析 & $d_i$、$D_{10}/D_{50}/D_{90}$、逐筛级占比 & 同批机械筛分逐筛称量真值 \\
形貌分类 & AR、圆度、凸度及投影形貌类别 & 人工/专家二维形貌标签；三维针片状需物理量规 \\
异常检测 & $>50$~mm、过小和疑似泥团候选 & 已知异常目标或人工语义标签 \\
现场约束 & 标定、推理、质检与报告导出 & 完整 SOP wall-clock time \\
\bottomrule
\end{tabular}
\end{table}

围绕上述问题，本文的贡献集中在三点。其一，复用 ImageGrains 2.0\cite{mair2026imagegrains2} 与 Cellpose-SAM\cite{pachitariu2025cellposeSAM} 的预训练颗粒实例感知能力，降低比赛现场重新标注和训练的依赖，并把模型输出限制为可检查的实例轮廓与几何量。其二，不把二维短轴直接视为标准筛孔粒径，而是提出低维、可解释、可由机械筛分真值校准的“筛分等效粒径 + 质量加权”映射，将视觉数量统计转换到标准筛分关注的质量级配。其三，将粒径、投影形貌、异常筛查和现场操作组织为同一条可追溯数据链，同时明确区分功能演示、软件一致性和物理精度三种证据等级，避免把未配对机械筛分的视觉结果过度解释为准确率。

下一节说明为什么选择“强预训练视觉基座 + 显式物理/统计校准”，以及各模块如何组成在线推理链和离线校准环。
